\documentclass[a4paper,12pt]{report}
\usepackage[utf8]{inputenc}
\usepackage[english, italian]{babel}
\usepackage{graphicx}
\usepackage{float}
\usepackage{tabularx}
\usepackage{makecell}
\usepackage{titlesec}
\usepackage{fancyhdr}
\usepackage{lastpage}
\usepackage{xurl}
\usepackage{hyperref}
\usepackage{geometry}
\usepackage[table,dvipsnames]{xcolor}

\setcounter{secnumdepth}{5}% Number \subsubsection
\setcounter{tocdepth}{5}

\titleformat{\chapter}[display]
{\normalfont\bfseries}{}{0pt}{\LARGE} %QUESTO SERVE PER AVERE SOLO IL NOME DEL CAPITOLO CHE VOGLIO IO

\titlespacing*{\chapter}{0cm}{0cm}{0.2cm}

\geometry{
	left=20mm,
	right=20mm,
}

\renewcommand{\contentsname}{Indice} %QUESTO SERVE PER AVERE L'INDICE CON IL NOME CHE VOGLIO IO

\pagestyle{plain}

\fancypagestyle{plain}{
	\fancyhf{}
	\lhead{\includegraphics[width=3cm]{../immagini/minilogo.jpg}}
	\chead{}
	\rhead{\fontsize{12}{10}Studio di fattibilità}
	\lfoot{}
	\cfoot{\thepage\ di \pageref{LastPage}}
	\rfoot{}
}

\pagestyle{plain}

\begin{document}
	\makeatletter
	\begin{titlepage}
		\begin{center}
			\vspace*{-4,0cm}
			\author{Jawa Druids}
			\title{Norme di Progetto}
			\date{}
			\includegraphics[width=0.7\linewidth]{../immagini/DRUIDSLOGO.jpg}\\[4ex]
			{\huge \bfseries  \@title }\\[2ex] 
			{\LARGE  \@author}\\[50ex]
			\vspace*{-8,0cm}
			\begin{table}[H]
				\centering
				\begin{tabular}{c|c}
					\textbf{Versione} & x.x.x \\
					\textbf{Data approvazione} & xx-xx-xxxx\\
					\textbf{Responsabile} & Nome Cognome\\
					\textbf{Redattori} & Andrea Cecchin \\
					\textbf{Verificatori} & \makecell{Nome Cognome \\ Nome Cognome} \\
					\textbf{Stato} & stato\\
					\textbf{Lista distribuzione} & \makecell{Jawa Druids \\ Tullio \\ GGGG}\\
					\textbf{Uso} & Interno			
				\end{tabular}
			\end{table}
			\fontsize{17}{12}\textbf{Sommario} \\
			Il documento redatto riferisce le regole, gli strumenti e le convenzioni a cui il gruppo JawaDruids ha stabilito di attenersi e seguire per l'intera durata dello sviluppo del progetto.
		\end{center}
	\end{titlepage}
	\makeatother
	\quad
\begin{center}
	\LARGE\textbf{Registro delle modifiche}
\end{center}
\def\tabularxcolumn#1{m{#1}}
{\rowcolors{2}{RawSienna!90!RawSienna!20}{RawSienna!70!RawSienna!40}

\begin{center}
	\renewcommand{\arraystretch}{1.4}
	\begin{longtable}{|p{4cm}|p{3cm}|p{2.8cm}|p{2cm}|p{1.7cm}|}
	\hline
	\rowcolor{airforceblue}
	\textbf{Modifica} & \textbf{Autore} & \textbf{Ruolo} & \textbf{Data} & \textbf{Versione}
	\\
	\hline
	\textit{Approvazione del documento per RR} & Igli Mezini & \textit{Responsabile} & 10-01-2021 & v1.0.0
	\\
	\hline
	\textit{Revisione Finale del documento} & Margherita Mitillo & \textit{Verificatore} & 09-01-2021 & v0.9.0
	\\
	\hline
	\textit{Verificato capitolo \S~\ref{Formazione}} & Margherita Mitillo & \textit{Verificatore} & 08-01-2021 & v0.8.0
	\\
	\hline
	\textit{Verificate sezioni \S~\ref{ProcessiPrimariProgettazione}, \S~\ref{ProcessiPrimariCodifica}, \S~\ref{ProcessiPrimariStrumenti}, \S~\ref{Standard ISO/IEC 15504} } & Andrea Dorigo & \textit{Verificatore} & 08-01-2021 & v0.7.0
	\\
	\hline
	\textit{Verificate sezioni \S~\ref{ProcessiDiSupportoGestioneDellaConfigurazione}, \S~\ref{ProcessiOrganizzativiProcessoDiPianificazione}, \S~\ref{ProcessiOrganizzativiFormazione}} & Margherita Mitillo & \textit{Verificatore} & 07-01-2021 & v0.6.0
	\\
	\hline
	\textit{Verificate sezioni \S~\ref{ProcessiDiSupportoVerifica}, \S~\ref{ProcessiDiSupportoValidazione}, \S~\ref{StandardISO/IEC9126}} & Emma Roveroni & \textit{Verificatore} & 07-01-2021 & v0.5.0
	\\
	\hline
	\textit{Aggiunto capitolo \S~\ref{Formazione}} & Mattia Cocco & \textit{Amministratore} & 06-01-2021 & v0.4.7
	\\
	\hline
	\textit{Aggiunte sezioni \S~\ref{ProcessiDiSupportoVerifica} e \S~\ref{ProcessiDiSupportoValidazione} } & Alfredo Graziano & \textit{Amministratore} & 05-01-2021 & v0.4.6
	\\
	\hline
	\textit{Aggiunto capitolo \S~\ref{ProcessiOrganizzativiFormazione}} & Mattia Cocco & \textit{Amministratore} & 05-01-2021 & v0.4.5
	\\
	\hline
	\textit{Aggiunte sezioni \S~\ref{ProcessiPrimariProgettazione}, \S~\ref{ProcessiPrimariCodifica}, \S~\ref{ProcessiPrimariStrumenti}} & Andrea Cecchin & \textit{Amministratore} & 03-01-2021 & v0.4.4
	\\
	\hline
	\textit{Aggiunta sezione \S~\ref{ProcessiDiSupportoGestioneDellaConfigurazione}} & Andrea Dorigo & \textit{Amministratore} & 02-01-2021 & v0.4.3
	\\
	\hline
	\textit{Aggiunte sezioni \S~\ref{ProcessiPrimariProspettiveAnalisiDeiRequisitiMetriche}, \S~\ref{ProcessiDiSupportoVerificaDescrizione}} & Mattia Cocco & \textit{Amministratore} & 02-01-2021 & v0.4.2
	\\
	\hline
	\textit{Aggiunto capitolo \S~\ref{Standard ISO/IEC 15504}} & Alfredo Graziano & \textit{Amministratore} & 01-01-2021 & v0.4.1
	\\
	\hline
	\textit{Verificate sezioni \S~\ref{ProcessiDiSupportoGestioneDellaQualità} e \S~\ref{ProcessiOrganizzativiProcessoDiCoordinamento} } & Emma Roveroni & \textit{Verificatore} & 31-12-2020 & v0.4.0
	\\
	\hline
	\textit{Aggiunta sezione \S~\ref{ProcessiOrganizzativiProcessoDiPianificazioneMetriche}} & Alfredo Graziano & \textit{Amministratore} & 30-12-2020 & v0.3.4
	\\
	\hline
	\textit{Aggiunte sezioni \S~\ref{ProcessiDiSupportoGestioneDellaQualità}, \S~\ref{StandardISO/IEC9126} } & Mattia Cocco & \textit{Amministratore} & 30-12-2020 & v0.3.3
	\\
	\hline
	\textit{Aggiunte sezioni \S~\ref{ProcessiOrganizzativiProcessoDiPianificazioneScopo}, \S~\ref{ProcessiOrganizzativiProcessoDiPianificazioneRuoliDiProgetto}, \S~\ref{ProcessiOrganizzativiProcessoDiPianificazioneAssegnazioneDeiCompiti}, \S~\ref{ProcessiOrganizzativiProcessoDiPianificazioneTrelloEGitkraken}, \S~\ref{ProcessiOrganizzativiProcessoDiPianificazioneStrumenti} } & Igli Mezini & \textit{Amministratore} & 30-12-2020 & v0.3.2
	\\
	\hline
	\textit{Aggiunte sezioni \S~\ref{ProcessiOrganizzativiProcessoDiCoordinamentoScopo}, \S~\ref{ProcessiOrganizzativiProcessoDiCoordinamentoComunicazione}, \S~\ref{ProcessiOrganizzativiProcessoDiCoordinamentoRiunioni}, \S~\ref{ProcessiOrganizzativiProcessoDiCoordinamentoStrumentiUtilizzatiPerIlProcessoDiCoordinamento} } & Igli Mezini & \textit{Amministratore} & 29-12-2020 & v0.3.1
	\\
	\hline
	\textit{Verificata sezione \S~\ref{ProcessiDiSupportoDocumentazione}} & Andrea Dorigo & \textit{Verificatore} & 15-12-2020 & v0.3.0
	\\
	\hline
	\textit{Verificate sezioni \S~\ref{ProcessiPrimariFornitura} e \S~\ref{ProcessiPrimariSviluppo} } & Margherita Mitillo & \textit{Verificatore} & 11-12-2020 & v0.2.0
	\\
	\hline
	\textit{Verificato capitolo \S~\ref{Introduzione}} & Margherita Mitillo & \textit{Verificatore} & 10-12-2020 & v0.1.0
	\\
	\hline
	\textit{Aggiunte sezioni \S~\ref{ProcessiDiSupportoDocumentazioneMetricheCorrettezzaOrtografica}, \S~\ref{ProcessiDiSupportoDocumentazioneDirectoryDiUnDocumento}} & Andrea Dorigo & \textit{Amministratore} & 08-12-2020 & v0.0.9
	\\
	\hline
	\textit{Aggiunte sezioni \S~\ref{ProcessiDiSupportoDocumentazioneMetriche}, \S~\ref{ProcessiDiSupportoDocumentazioneStrumentiDiStesura}} & Igli Mezini & \textit{Amministratore} & 03-12-2020 & v0.0.8
	\\
	\hline
	\textit{Aggiunte sezioni \S~\ref{ProcessiDiSupportoDocumentazioneStrutturaGeneraleDeiDocumenti}, \S~\ref{ProcessiDiSupportoDocumentazioneNormeTipografiche}, \S~\ref{ProcessiDiSupportoDocumentazioneElementiGrafici}} & Igli Mezini & \textit{Amministratore} & 01-12-2020 & v0.0.7
	\\
	\hline
	\textit{Aggiunta sezione \S~\ref{ProcessiPrimariSviluppo}} & Andrea Cecchin & \textit{Amministratore} & 30-11-2020 & v0.0.6
	\\
	\hline
	\textit{Aggiunte sezioni \S~\ref{ProcessiDiSupportoDocumentazioneTemplateInFormatoLatex}, \S~\ref{ProcessiDiSupportoDocumentazioneDocumentiProdotti}, \S~\ref{ProcessiDiSupportoDocumentazioneDirectoryDiUnDocumento}} & Igli Mezini & \textit{Amministratore} & 29-11-2020 & v0.0.5
	\\
	\hline
	\textit{Aggiunte sezioni \S~\ref{ProcessiDiSupportoDocumentazioneDescrizione}, \S~\ref{ProcessiDiSupportoDocumentazioneImplementazioneDelDocumento}, \S~\ref{ProcessiDiSupportoDocumentazioneCicloDiVitaDiUnDocumento}} & Igli Mezini & \textit{Amministratore} & 28-11-2020 & v0.0.4
	\\
	\hline
	\textit{Aggiunte sezioni \S~\ref{ProcessiPrimariFornituraStudioDiFattibilità}, \S~\ref{ProcessiPrimariFornituraAltraDocumentazioneDaFornire}, \S~\ref{ProcessiPrimariFornituraStrumenti}} & Andrea Cecchin & \textit{Amministratore} & 26-11-2020 & v0.0.3
	\\
	\hline
	\textit{Aggiunte sezioni \S~\ref{Introduzione} e \S~\ref{ProcessiPrimariFornituraScopo}} & Andrea Cecchin & \textit{Amministratore} & 25-11-2020 & v0.0.2
	\\
	\hline
	\textit{Prima stesura del documento} & Andrea Cecchin & \textit{Amministratore} & 24-11-2020 & v0.0.1
	\\
	\hline
	\end{longtable}
\end{center}

	
	
	\tableofcontents{}
	
	\chapter{Introduzione}\label{Introduzione}

\section{Scopo del documento}
Il Piano di Qualifica è un documento su cui si prevede di operare per l’intera durata del progetto 
e il cui scopo è presentare e descrivere le strategie di verifica e validazione adottate 
dal gruppo Jawa Druids al fine di garantire la qualità di prodotto e di processo.  
Per raggiungere questo obbiettivo viene applicato un sistema di verifica continua sui processi in corso e 
sulle attività svolte, in modo da rilevare e correggere subito eventuali anomalie, riducendo lo spreco di risorse 
ed il rischio di reiterare gli stessi errori.
\section{Scopo del prodotto}
L'obiettivo del prodotto \textit{GDP - Gathering Detection Platform} di \textit{Sync Lab} è la creazione di un sistema software in 
grado di acquisire, monitorare, utilizzare e correlare tra loro tutti i dati acquisiti tramite sensoristica (come telecamere di videosorveglianza)
e altre sorgenti dati (es. orari con il maggior numero di visite degli esercizi commerciali estrapolati da Google) con lo
scopo di avere indicazioni di potenziali assembramenti e quindi di fornire supporto per trovare soluzioni che limitino tali condizioni.
\section{Glossario}
Al fine di favorire la comprensione totale della documentazione relativa al progetto, viene fornito un documento denominato \textit{Glossario -- fornire versione}
che riporta la spiegazione della terminologia specifica e tecnica utilizzata per la stesura di questo ed altri documenti.
Tali termini, quando utilizzati, vengono contrassegnati con un apice \textsuperscript{G} alla fine della parola.
\section{Standard di progetto}
Decidere se scrivere di aver utilizzato un determinato standard ISO per il software, se si, definirlo nelle norme di progetto!
\section{Riferimenti}
\subsection{Riferimenti normativi}
\begin{itemize}
	\item \textit{Norme di Progetto}
\end{itemize}
\subsection{Riferimenti informativi}
\begin{itemize}
	\item \textit{IEEE Recommended Practice for Software Requirements Specifications:}\\
	\url{https://ieeexplore.ieee.org/document/720574}
	\item \textit{Seminario per approfondimenti tecnici del capitolato C3:}\\
	\url{https://www.math.unipd.it/~tullio/IS-1/2020/Progetto/ST1.pdf}		
\end{itemize}
	\chapter{Processi Primari}\label{ProcessiPrimari}
\section{Fornitura}\label{ProcessiPrimariFornitura}
\subsection{Scopo}\label{ProcessiPrimariFornituraScopo}
La fornitura secondo lo standard ISO/IEC 12207:1995 descrive le attività$_G$ e i compiti svolti dal fornitore al fine di sviluppare un prodotto soddisfacente e che rispetti appieno le richieste del committente$_G$.
Durante questa fase si prevede la compilazione di diversi documenti, i quali verranno inviati al committente$_{\scaleto{G}{3pt}}$ per guadagnare la possibilità di lavorare al progetto offerto dall'azienda \textit{Sync Lab}.
Il fornitore esegue un'attività$_{\scaleto{G}{3pt}}$ di analisi e stesura dello \textit{Studio di Fattibilità}, documento che rileva i rischi e le criticità riscontrate nella richiesta di appalto.
Si definisce inoltre un accordo contrattuale con il proponente$_G$ mediante il quale si regolano i rapporti con l'azienda, la consegna e la manutenzione del prodotto sviluppato. 

\subsection{Studio di Fattibilità}\label{ProcessiPrimariFornituraStudioDiFattibilità}
Lo \textit{Studio di Fattibilità} consiste nell'analisi e nella valutazione sistematica delle caratteristiche, dei costi, e dei possibili risultati di un progetto sulla base di una preliminare idea di massima.
A seguito della presentazione dei capitolati d'appalto da parte di ogni proponente$_{\scaleto{G}{3pt}}$ avvenuta il 05-11-2020, il \textit{Responsabile di Progetto} si è impegnato a programmare incontri con tutti i componenti del gruppo \textit{Jawa Druids} per valutare le scelte di ogni membro e attuare così un primo scambio di idee. Una volta individuato il capitolato d'interesse gli \textit{Analisti} hanno provveduto alla stesura dello \textit{Studio di Fattibilità}, i quali hanno fornito un'analisi accurata dei capitolati presentati.
Nella stesura dello \textit{Studio di Fattibilità} per ogni capitolato$_{\scaleto{G}{3pt}}$ si riporterà:
\begin{itemize}
	\item informazioni generali: informazioni riguardanti il proponente$_{\scaleto{G}{3pt}}$;
	\item descrizione del capitolato$_{\scaleto{G}{3pt}}$: sintesi del progetto da sviluppare; 
	\item finalità del progetto: finalità richieste dal capitolato$_{\scaleto{G}{3pt}}$ d'appalto;
	\item tecnologie interessate: tecnologie che verranno utilizzate nello svolgimento del capitolato$_{\scaleto{G}{3pt}}$
	\item aspetti positivi: aspetti favorevoli alla scelta del capitolato$_{\scaleto{G}{3pt}}$;
	\item criticità e fattori di rischio: problematiche che potrebbero sorgere durante lo svolgimento del capitolato;
	\item conclusioni: accettazione o rifiuto del capitolato$_{\scaleto{G}{3pt}}$ in base alle informazioni illustrate precedentemente e anche all'interesse dimostrato da ogni membro nel gruppo.
\end{itemize}
\subsection{Altra documentazione da fornire}\label{ProcessiPrimariFornituraAltraDocumentazioneDaFornire}
Oltre allo \textit{Studio di Fattibilità} vengono consegnati altri documenti all'azienda \textit{Sync Lab} ed ai committenti \textit{Prof. Tullio Vardanega} e \textit{Prof. Riccardo Cardin}. Questi documenti sono necessari al fine di tracciare le attività$_{\scaleto{G}{3pt}}$ di Analisi, Pianificazione, Verifica, Validazione e Controllo di Qualità per assicurare una completa trasparenza durante tutta la durata del ciclo di vita del progetto.
I documenti in questione sono:
\begin{itemize}
	\item \textit{Analisi dei Requisiti}: identifica e dettaglia in modo completo ed esaustivo i requisiti$_G$ del sistema descritto nel capitolato che il fornitore si impegna a soddisfare;
	\item \textit{Piano di Qualifica}: illustra la strategia complessiva di verifica e validazione proposta dal fornitore per pervenire al collaudo del sistema con la massima efficienza ed efficacia;
	\item \textit{Piano di Progetto}: presenta l'organigramma dettagliato del fornitore, lo schema proposto per l'assegnazione e la rotazione dei ruoli di progetto, l'impegno complessivo previsto per ogni ruolo e per ogni individuo, l'analisi dei rischi, la pianificazione di massima per la realizzazione del prodotto, e il corrispondente conto economico preventivo.
\end{itemize} 
Alla documentazione appena illustrata il gruppo \textit{Jawa Druids} allegherà inoltre una lettera di presentazione con la quale si formalizza l'impegno nel portare al termine il capitolato prescelto entro i termini definiti nella lettera e rispettandone i requisiti$_{\scaleto{G}{3pt}}$ minimi.
\subsection{Strumenti}\label{ProcessiPrimariFornituraStrumenti}
Di seguito sono riportati gli strumenti impiegati dal gruppo durante il progetto per il processo di fornitura.
\subsubsection{Documenti Google}\label{ProcessiPrimariFornituraStrumentiDocumentiGoogle}
Questo strumento viene utilizzato per la realizzazione di documenti in cui più persone possono lavorarci contemporaneamente osservando le modifiche in tempo reale. 
\begin{center}
	\url{https://docs.google.com}
\end{center}
\subsubsection{Fogli Google}\label{ProcessiPrimariFornituraStrumentiFogliGoogle}
Questo strumento viene utilizzato per la realizzazione di grafici e fogli di calcolo in cui più persone possono lavorarci contemporaneamente osservando le modifiche in tempo reale.
\begin{center}
	\url{https://www.google.com/sheets}
\end{center}


\section{Sviluppo}\label{ProcessiPrimariSviluppo}
\subsection{Scopo}\label{ProcessiPrimariScopo}
Il processo di sviluppo contiene tutte le attività$_{\scaleto{G}{3pt}}$ che riguardano la produzione del software richiesto dal cliente, in particolare analisi dei requisiti$_{\scaleto{G}{3pt}}$, design, codifica, integrazione, test e installazione.
\subsection{Descrizione}\label{ProcessiPrimariDescrizione}
Di seguito vengono elencate le varie attività$_{\scaleto{G}{3pt}}$ che caratterizzano tale processo:
\begin{itemize}
	\item Analisi dei requisiti$_{\scaleto{G}{3pt}}$;
	\item Progettazione architetturale;
	\item Codifica del software.
\end{itemize}
\subsection{Prospettive}\label{ProcessiPrimariProspettive}
Le prospettive alla fine della stesura del processo in questione sono le seguenti:
\begin{itemize}
	\item individuare e stabilire gli obbiettivi di sviluppo;
	\item individuare e stabilire i vincoli tecnologici;
	\item individuare e stabilire i vincoli di design;
	\item produrre un prodotto finale che rispecchi gli obiettivi imposti nello sviluppo e che superi i test e i controlli di qualità stabiliti dal proponente$_{\scaleto{G}{3pt}}$.
\end{itemize}
\subsubsection{Analisi dei requisiti}\label{ProcessiPrimariProspettiveAnalisiDeiRequisiti}
\paragraph{Scopo}\label{ProcessiPrimariProspettiveAnalisiDeiRequisitiScopo}\mbox{}\\
L'analisi dei requisiti$_{\scaleto{G}{3pt}}$ viene redatto dagli \textit{Analisti}, lo scopo è quello di definire le funzionalità che il nuovo prodotto deve offrire, ovvero i requisiti$_{\scaleto{G}{3pt}}$ che devono essere soddisfatti dal software sviluppato.\\
Gli obiettivi della stesura dell'\textit{Analisi dei Requisiti} sono:
\begin{itemize}
	\item stabilire lo scopo nello sviluppo del prodotto;
	\item definire riferimenti precisi ed affidabili ai \textit{Progettisti};
	\item stabilire i requisiti$_{\scaleto{G}{3pt}}$ e le funzionalità concordate con il cliente;
	\item individuare per i \textit{Verificatori} riferimenti per le attività$_{\scaleto{G}{3pt}}$ di controllo dei test.
\end{itemize}
\paragraph{Descrizione}\label{ProcessiPrimariProspettiveAnalisiDeiRequisitiDescrizione}\mbox{}\\
I requisiti$_{\scaleto{G}{3pt}}$ possono essere individuati in diverse fonti, quali:
\begin{itemize}
	\item \textit{Capitolato$_{\scaleto{G}{3pt}}$ d'Appalto}: i requisiti$_{\scaleto{G}{3pt}}$ sono stati individuati attraverso la lettura del documento fornito dal proponente$_{\scaleto{G}{3pt}}$ \textit{Sync Lab} sul capitolato proposto;
	\item \textit{Verbali Interni}: attraverso le riunioni attuate internamente dagli \textit{Analisti} sono emersi vari requisiti$_{\scaleto{G}{3pt}}$;
	\item \textit{Verbali Esterni}: attraverso contatti e discussioni effettuate con il responsabile aziendale Fabio Pallaro sono emersi requisiti$_{\scaleto{G}{3pt}}$, i quali vi sarà assegnato un codice presente nella tabella dei tracciamenti;
\end{itemize}
\paragraph{Prospettive}\label{ProcessiPrimariProspettiveAnalisiDeiRequisitiProspettive}\mbox{}\\
L'obiettivo dell'\textit{Analisi dei Requisti} è quello di redigere un documento che racchiuda al suo interno tutti i requisiti$_{\scaleto{G}{3pt}}$ richiesti dal proponente$_{\scaleto{G}{3pt}}$.
\paragraph{Struttura}\label{ProcessiPrimariProspettiveAnalisiDeiRequisitiStruttura}\mbox{}\\ %wip non è ancora perfettamente definito ???
L'\textit{Analisi dei requisiti$_{\scaleto{G}{3pt}}$} è strutturato nel seguente modo:
\begin{itemize}
	\item \textbf{Introduzione}: in questa sezione si introduce lo scopo del documento e riferimenti a documenti esterni;
	\item \textbf{Descrizione generale}: il prodotto viene descritto, si individuano le fasi generali del progetto e l'utenza a cui è destinato il prodotto;
	\item \textbf{Fasi del prodotto}: si elencano in gruppi le fasi del progetto, le quali vengono suddivise in parti per individuare più dettagliatamente tutti i processi di sviluppo del software;
	\item \textbf{Requisiti$_{\scaleto{G}{3pt}}$}: utilizzando le fasi descritte e tutti i documenti sopra indicati si dettagliano tutti i requisiti$_{\scaleto{G}{3pt}}$ obbligatori, facoltativi e desiderabili da implementare.
\end{itemize}
\paragraph{Classificazione dei requisiti}\label{ProcessiPrimariProspettiveAnalisiDeiRequisitiClassificazioneDeiRequisiti}\mbox{}\\% RS1.2 = rs(requisiti specifici) 1(fase generale a cui si riferiscono) 2(numero del requisito incrementale)
I requisiti$_{\scaleto{G}{3pt}}$ sono stati individuati utilizzando la seguente codifica:
\begin{center}
	\textbf{RS[classificazione][tipo\_di\_requisito][codice\_requisito]}
\end{center} 
La descrizione della classificazione è la seguente:
\begin{itemize}
	\item \textbf{RS}: è l'acronimo per Requisito$_{\scaleto{G}{3pt}}$ Specifico;
	\item \textbf{Classificazione}: individua la classificazione del requisito$_{\scaleto{G}{3pt}}$:
	\begin{itemize}
		\item Funzionale: indicato dalla lettera "F";
		\item Qualità: indicato dalla lettera "Q";
		\item Vincolo: indicato dalla lettera "V";
		\item Prestazionale: indicato dalla lettera "P".
	\end{itemize}
	\item \textbf{Tipo\_di\_requisito$_{\scaleto{G}{3pt}}$}: individua la tipologia di requisito$_{\scaleto{G}{3pt}}$:
	\begin{itemize}
		\item Obbligatorio: indicato con la lettera "O" individua un requisito$_{\scaleto{G}{3pt}}$ essenziale allo sviluppo del progetto e necessario al suo completamento;
		\item Desiderabile: indicato con la lettera "D" individua un requisito$_{\scaleto{G}{3pt}}$ utile al prodotto e che dà valore aggiunto ad esso, ma non essenziale al suo completamento;
		\item Facoltativo: indicato con la lettera "F" individua un requisito$_{\scaleto{G}{3pt}}$ che può essere sviluppato, ma può anche non essere completato;
	\end{itemize}
	\item \textbf{Codice\_requisito}: è rappresentato da un codice identificativo univoco nella forma gerarchica padre/figlio.
\end{itemize}
Ogni requisito$_{\scaleto{G}{3pt}}$ è strutturato nella tabella nel seguente modo:
\begin{itemize}
	\item \textbf{Identificativo}: individua univocamente il requisito$_{\scaleto{G}{3pt}}$;
	\item \textbf{Descrizione}: descrizione del requisito$_{\scaleto{G}{3pt}}$ richiesto nella determinata fase;
	\item \textbf{Tipo di requisito$_{\scaleto{G}{3pt}}$}: individua la tipologia di requisito$_{\scaleto{G}{3pt}}$ obbligatorio, desiderabile e facoltativo;
	\item \textbf{Fonte}: ogni requisito$_{\scaleto{G}{3pt}}$ è ricavato da una delle seguenti fonti:
	\begin{itemize}
		\item Capitolato$_{\scaleto{G}{3pt}}$: individua il documento del capitolato$_{\scaleto{G}{3pt}}$;
		\item Interno: il requisito$_{\scaleto{G}{3pt}}$ è stato individuato dagli analisti;
		\item Verbale: si tratta del documento in riferimento alla discussione con il proponente$_{\scaleto{G}{3pt}}$;
		\item Fasi del capitolato$_{\scaleto{G}{3pt}}$: il requisito$_{\scaleto{G}{3pt}}$ è stato individuato dalla fase del capitolato$_{\scaleto{G}{3pt}}$ individuata con il proprio codice.
	\end{itemize}
\end{itemize}
\paragraph{Classificazione delle fasi}\label{ProcessiPrimariProspettiveAnalisiDeiRequisitiClassificazioneDelleFasi}\mbox{}\\% FC1.2 = fc(fasi capitolato) 1 (macro-fase) 2(fase precisa)
La descrizione dettagliata delle fasi è stata incorporata nel documento per individuare la fase in cui il requisito$_{\scaleto{G}{3pt}}$ è necessario che venga implementato o sia possibile l'implementazione.\\
La classificazione delle fasi è definita nel seguente modo:
\begin{center}
	\textbf{FC[codice\_macro-fase].[codice\_fase]}
\end{center}
La descrizione della classificazione è la seguente:
\begin{itemize}
	\item \textbf{FC}: è l'acronimo per Fasi Capitolato$_{\scaleto{G}{3pt}}$;
	\item \textbf{Codice\_macro-fase}: individua un numero riferito alla fase dello sviluppo del progetto;
	\item \textbf{Codice\_fase}: individua un numero riferito ad una suddivisione della fase del progetto.
\end{itemize}
Ogni suddivisione delle fasi è strutturato nel seguente modo:
\begin{itemize}
	\item \textbf{Principali}:
	\begin{itemize}
		\item \textbf{Identificativo}: inserito prima del nome della fase per individuarla univocamente;
		\item \textbf{Nome}: identifica il nome della fase;
		\item \textbf{Descrizione}: descrizione dello scopo della fase in analisi.
	\end{itemize}
	\item \textbf{Aggiuntivi}:
\begin{itemize}
	\item \textbf{Input}: dati in ingresso necessari allo sviluppo della fase;
	\item \textbf{Output}: dati in uscita elaborati o modificati durante la fase;
	\item \textbf{Linguaggio di programmazione}: linguaggio di programmazione utilizzato per lo sviluppo della fase;
	\item \textbf{Processo}: sono le operazioni svolte nella fase;
	\item \textbf{Risposta ad errori}: operazione della fase nel caso di errori;
	\item \textbf{Tipi di algoritmi}: individua gli algoritmi utilizzabili;
	\item \textbf{Strumenti}: sono strumenti utilizzati nello svolgimento della fase;
	\item \textbf{Vincolo}: vincoli individuati dal gruppo o dal proponente$_{\scaleto{G}{3pt}}$.
\end{itemize}
\end{itemize}
\paragraph{Qualità dei requisiti}\label{ProcessiPrimariProspettiveAnalisiDeiRequisitiQualitàDeiRequisiti}\mbox{}\\
Ogni requisito$_{\scaleto{G}{3pt}}$ deve rispettare le seguenti qualità:
\begin{itemize}
	\item devono essere correttamente descritti;
	\item non devono essere ambigui, ogni requisito$_{\scaleto{G}{3pt}}$ fornirà un'unica interpretazione;
	\item devono essere completi, cioè descrivere in modo completo e in tutte le sue parti la funzionalità da implementare;
	\item ogni requisito$_{\scaleto{G}{3pt}}$ deve essere consistente, cioè non deve avere conflitti con altri requisiti$_{\scaleto{G}{3pt}}$ individuati;
	\item devono essere modificabili, cioè ogni requisito$_{\scaleto{G}{3pt}}$ nel corso dello sviluppo del progetto può essere rivalutato e modificato e bisogna mantenere uno storico dei cambiamenti;
	\item ogni requisito$_{\scaleto{G}{3pt}}$ deve essere tracciabile, cioè ogni requisito$_{\scaleto{G}{3pt}}$ deve essere tracciabile ad ogni suo test o codice di implementazione o alla sua origine;
	\item ogni requisito$_{\scaleto{G}{3pt}}$ deve essere classificato per importanza o stabilità;
	\item ogni requisito$_{\scaleto{G}{3pt}}$ deve essere verificabile, cioè deve essere possibile una sua verifica attraverso un processo nel quale una persona o una macchina può controllarlo.
\end{itemize}

\paragraph{Metriche}\label{ProcessiPrimariProspettiveAnalisiDeiRequisitiMetriche}\mbox{}\\
Con \textbf{MQPD01} si intende la \textbf{Totalità delle implementazioni}. Indice riportante l'interezza del prodotto software, rispetto ai requisiti$_{\scaleto{G}{3pt}}$ posti, mediante un valore percentuale.:

\begin{center}
	\textbf{T=($1-\frac{RnI}{RI}$)*100}
\end{center}
Dove:
\begin{itemize}
	\item \textbf{T} sta per \textit{Totalità}, riferito ai requisiti$_{\scaleto{G}{3pt}}$ da implementare;
	\item \textbf{RnI} sta per \textit{Requisito$_{\scaleto{G}{3pt}}$ non Implementato};
	\item \textbf{RI} sta per \textit{Requisito$_{\scaleto{G}{3pt}}$ Implementato}.
\end{itemize}
I range accettabili per il risultato di \textbf{T} sono così suddivisi:
\begin{itemize}
	\item 90\% $<$ \textbf{T} $\leq$ 100\% indica che la copertura dei requisiti$_{\scaleto{G}{3pt}}$ proposti è quasi totale;
	\item 80\% $<$ \textbf{T} $\leq$ 90\% indica che la copertura dei requisiti$_{\scaleto{G}{3pt}}$ proposti è sufficiente, buona;
	\item \textbf{T} $\leq$ 80\% indica che la copertura dei requisiti$_{\scaleto{G}{3pt}}$ proposti è insufficiente.
\end{itemize}


\subsection{Progettazione}\label{ProcessiPrimariProgettazione}
\subsubsection{Scopo}\label{ProcessiPrimariProgettazioneScopo}
La Progettazione è un'attività$_{\scaleto{G}{3pt}}$ svolta dai \textit{Progettisti}. In questa fase si individuano, attraverso l'\textit{Analisi dei Requisiti}, le caratteristiche che il prodotto deve avere per soddisfare tutti i requisiti$_{\scaleto{G}{3pt}}$ richiesti dal proponente$_{\scaleto{G}{3pt}}$. Lo scopo è quello di determinare la soluzione migliore per risolvere ogni requisito$_{\scaleto{G}{3pt}}$ individuato. 
\subsubsection{Aspettative}\label{ProcessiPrimariProgettazioneAspettative}
Alla conclusione della stesura della progettazione si ha come risultato la realizzazione dell'architettura del software da sviluppare. Tale architettura è necessaria ai programmatori per individuare le istruzioni necessarie a sviluppare il prodotto finito.
\subsubsection{Descrizione}\label{ProcessiPrimariProgettazioneDescrizione}
La progettazione è suddivisa principalmente in due parti:
\begin{itemize}
	\item \textbf{Technology baseline$_G$}: contiene tutte le specifiche adottate nella progettazione ad alto livello e delle sue componenti, i test di verifica e l'elenco dei diagrammi UML$_G$ utilizzati per la definizione dell'architettura;
	\item \textbf{Product baseline$_G$}: perfeziona ulteriormente la fase di progettazione, integrando ciò che è riportato nella technology baseline$_{\scaleto{G}{3pt}}$, definendo inoltre i test necessari alla verifica.
\end{itemize}

\subsubsection{Technology baseline}\label{ProcessiPrimariProgettazioneTecnologyBaseline}
La technology baseline$_{\scaleto{G}{3pt}}$ viene redatta dal \textit{Progettista} ed include:
\begin{itemize}
	\item \textbf{Diagramma UML$_{\scaleto{G}{3pt}}$}: diagrammi sulle classi usate, package, attività$_{\scaleto{G}{3pt}}$ e sequenze;
	\item \textbf{Tecnologie}: descrizione delle tecnologie utilizzate nel progetto, i loro vantaggi e svantaggi;
	\item \textbf{Design pattern}: descrizione di ogni design pattern per la realizzazione dell'architettura, ognuno di essi viene accompagnato da un diagramma che ne espone il significato e la struttura; 
	\item \textbf{Tracciamento delle componenti}: ogni componente viene riferito al requisito$_{\scaleto{G}{3pt}}$ che soddisfa;
	\item \textbf{Test di integrazione}: corrispondono a tutti i test necessari alla verifica del sistema in tutte le sue unità.
\end{itemize}
\subsubsection{Product baseline}\label{ProcessiPrimariProgettazioneProductBaseline}
La product baseline$_{\scaleto{G}{3pt}}$ viene redatta dal \textit{Progettista} ed include:
\begin{itemize}
	\item \textbf{definizioni delle classi e funzioni}: descrizione univoca delle classi e funzioni evitando ridondanze;
	\item \textbf{tracciamento}: ogni requisito$_{\scaleto{G}{3pt}}$ viene tracciato attraverso le classi e funzioni, questo è necessario per individuare se è stato soddisfatto il requisito$_{\scaleto{G}{3pt}}$;
	\item \textbf{test di unità}: definiscono test per la verifica delle funzionalità del prodotto.
\end{itemize}
\subsection{Codifica}\label{ProcessiPrimariCodifica}

\subsubsection{Scopo}\label{ProcessiPrimariCodificaScopo}
La fase di codifica è la scrittura del codice per sviluppare la miglior soluzione del prodotto. In questa sezione si introdurranno tutte le norme necessarie allo sviluppo di un codice uniformato tra tutti i \textit{Programmatori} e rispettoso delle regole standard indicate nel documento.
\subsubsection{Aspettative}\label{ProcessiPrimariCodificaAspettative}
Conclusa la fase di codifica ci si attende un codice pulito e facile da leggere, utile nelle successive validazioni, modifiche e per agevolare la sua manutenzione. L'obiettivo è quello di sviluppare un prodotto conforme alle richieste individuate con il proponente$_{\scaleto{G}{3pt}}$. 

\subsubsection{Descrizione}\label{ProcessiPrimariCodificaDescrizione}
In questa fase la programmazione del prodotto dovrà rispettare le norme descritte nel documento. Perseguendo gli obiettivi individuati nel documento \textit{Piano di qualifica v1.0.0} si produrrà un software con un'alta qualità di codice.

\subsubsection{Stile di codifica}\label{ProcessiPrimariCodificaStileDiCodifica}
All'interno di questo paragrafo vengono elencate le norme da rispettare da ogni membro del gruppo per raggiungere uniformità del codice:

\begin{itemize}
	\item \textbf{Indentazione}: ogni blocco di codice scritto per il prodotto da sviluppare deve essere ben indentato e deve rispettare una misura di 4 spazi (equivalenti alla pressione del tasto \textit{tab}). Fanno eccezione a questa regola i commenti che vengono inseriti per spiegazioni di blocchi di codice;
	\item \textbf{Parentesizzazione}: le parentesi devono essere utilizzate in linea col blocco di codice scritto e non in una linea sottostante per separarlo;
	\item \textbf{Univocità delle variabili, metodi e funzioni}: ogni variabile, metodo e funzione utilizzata deve avere un nome significativo, esplicativo ed univoco;
	\item \textbf{Classi}: ogni classe deve avere il proprio nome scritto con l'iniziale maiuscola;
	\item \textbf{Metodi e funzioni}: il nome di metodi e funzioni devono iniziare per lettera minuscola e se composti da più parole le successive devono essere scritte con lettera maiuscola (stile \textit{CamelCase});
	\item \textbf{Lingua}: i commenti devono essere scritti in lingua inglese.
\end{itemize}

\subsection{Strumenti}\label{ProcessiPrimariStrumenti}

\subsubsection{PragmaDB}\label{ProcessiPrimariStrumentiPragmaDB}
Programma utilizzato per il tracciamento dei requisiti$_{\scaleto{G}{3pt}}$.
\begin{center}
	\url{https://pragmadb.com/}
\end{center}
\subsubsection{StarUML}\label{ProcessiPrimariStrumentiDrawIo}
Questo strumento viene usato per la realizzazione di diagrammi UML$_{\scaleto{G}{3pt}}$ in quanto è stato ritenuto semplice da utilizzare.
\begin{center}
	\url{https://staruml.io/}
\end{center}
\subsubsection{Atom}\label{ProcessiPrimariStrumentiAtom}
Questo strumento viene usato per la codifica del linguaggio Java e Javascript. Offre la piena compatibilità con Linux, Windows, macOS e fornisce molte integrazioni aggiuntive.
\begin{center}
	\url{https://atom.io}
\end{center}
	
	% bibliography, glossary and index would go here.
	
\end{document}